% 统计学习中的两层随机性：训练集决定假设，未来样本决定单次损失。
\begin{tikzpicture}[
  main node/.style={book node,minimum width=28mm,minimum height=11mm,inner sep=4pt},
  small node/.style={book node,fill=BookPaper,minimum width=24mm,minimum height=9mm,inner sep=3pt},
  note/.style={font=\kaishu\scriptsize,text=BookMuted,align=center}
]
  \node[main node,fill=FigureInput!12] (dist) {未知分布$\mathcal P$};
  \node[small node,below left=14mm and 9mm of dist] (sample) {训练集\\$S\sim\mathcal P^n$};
  \node[small node,below right=14mm and 9mm of dist] (future) {未来观测\\$z\sim\mathcal P$};
  \node[main node,fill=FigureModel!12,below=15mm of sample] (alg) {算法$A$};
  \node[main node,fill=FigureOutput!8,right=14mm of alg] (hyp) {输出$\hat h=A(S)$};
  \node[main node,fill=FigureLoss!12,below=13mm of hyp] (loss) {损失$\ell(\hat h,z)$};

  \draw[book arrow] (dist) -- node[note,left]{重复抽取\\训练数据} (sample);
  \draw[book arrow] (dist) -- node[note,right]{部署时的\\新数据} (future);
  \draw[book arrow] (sample) -- (alg);
  \draw[book arrow] (alg) -- (hyp);
  \draw[book arrow] (hyp) -- (loss);
  \draw[book arrow] (future) -- (loss);

  \draw[BookBlue,dashed,rounded corners]
    ([xshift=-4mm,yshift=4mm]sample.north west)
    rectangle ([xshift=4mm,yshift=-4mm]hyp.south east);
  \node[note,text=BookBlue,below=2mm of alg] {训练集随机性：决定最终得到哪个假设};
  \node[note,text=BookAmber,below=2mm of loss] {未来样本随机性：决定这一次预测是否出错};
\end{tikzpicture}
